\documentclass[12pt]{article}
\usepackage[utf8]{inputenc}
\usepackage[T1]{fontenc}
\usepackage{hyperref}
\title{Waneci Language Resources for K-12 Curriculum Development}
\author{Compiled by Claude Sonnet 5.5 (claude-sonnet-5-5) for Andr\'as Kornai}
\date{}
\begin{document}
\maketitle

\section{Introduction}
Waneci (also Wanetsi, Tarino, Tsalgari, Chalgari, Vanechi; ISO 639-3 wne) is an Eastern Iranian language of the Indo-European family. It is spoken mainly in northern Balochistan, Pakistan, in Harnai District and the Chawter area of Sanjawi. Estimates run from about 141,000 (Omniglot) to 150,000 (Wikipedia). Some scholars treat it as a Pashto dialect and others as a separate language. The first known linguistic research was by Georg Morgenstierne in 1929. One source describes it as at risk because of limited attention, and UNESCO's atlas lists it with a vulnerability rating.

\section{Primary Educational Resources}
I found no primers, textbooks or school programmes for Waneci.

\section{Linguistic and Reference Materials}
Overview: \url{https://en.wikipedia.org/wiki/Wanetsi}. Omniglot has a page on Wanetsi with its writing (in a Perso-Arabic style) and background: \url{https://www.omniglot.com/writing/wanetsi.htm}. The ASJP database has a wordlist entry: \url{https://asjp.clld.org/languages/WANECI}. Glottolog lists a reference by Hallberg (1992): \url{https://glottolog.org/resource/reference/id/120876}. A comparative database entry is at \url{https://lemontree210.pythonanywhere.com/en/doculect/waneci}.

\section{Religious and Community Resources}
No scripture translation or community organisation was found in the search results.

\section{Digital and Audio Resources}
The UNESCO World Atlas of Languages has a record: \url{https://ar.wal.unesco.org/node/73370}. I saw no audio, corpus or software resources.

\section{Curriculum Development Context}
Documentation is thin. Pashto and Urdu dominate schooling in the region. A curriculum would first need community-led standardization of an orthography and basic documentation (wordlists, texts, recordings). Whether Waneci is treated as a Pashto variety may affect planning and should be settled with the community.

\end{document}